\paragraph{\trGrading{}与符号。}\label{sub:grading-and-signs}
我们仍处在\cref{thm:fixed-point-set-assumption}的情形之下，但额外要求 $\dim(M^{\mathrm{fix}})=\dim(M)-2>0$，以避免平凡的例外。设 $M$ 带有一个（不必 $\iota_M$-不变的）二次复\trVolume{}形式 $\eta_M^2$。由于\trNormalBundle{}的平凡性，这会在\trFixedPoint{}集上诱导出一个类似的形式；用局部平方根的语言来说，
\begin{equation*}
\eta_M|M^{\mathrm{fix}}=\pm\pi^*(dz)\wedge\eta_{M^{\mathrm{fix}}}.
\end{equation*}
设 $L$ 是如\cref{thm:normal-angle-assumption}中的那样一个 $\iota_M$-不变的\trExact\trLagrange\trSubmanifold，且与\trFixedPoint{}集相交。$L$ 上的任何\trBraneStructure{}$(\alpha^{\#},P^{\#})$ 都在\trFixedPoint{}集 $L^{\mathrm{fix}}\subset M^{\mathrm{fix}}$ 上诱导出一个\trBraneStructure{}：\trGrading{}由 $\alpha^{\mathrm{fix},\#}=\alpha^{\#}|L^{\mathrm{fix}}-\frac{1}{\pi}\sigma$ 给出，其中 $\sigma$ 是\trNormalAngle{}。至于\trPinStructure，我们取与\trTangentBundle{}相伴的\trFrame{}丛 $P\to L$ 与 $P^{\mathrm{fix}}\to L^{\mathrm{fix}}$，并把第二个\trEmbedding{}到第一个之中：把一个 $L^{\mathrm{fix}}$-\trFrame{}$(X_1,\ldots,X_{n-1})$ 送到 $(e^{j\sigma},X_1,\ldots,X_{n-1})$，其中第一个分量是法方向上的一个向量。\trDoubleCover{}$P^{\#}\to P$ 沿这一\trEmbedding{}的\trPullback{}，连同它所继承的 $Pin_{n-1}\subset Pin_n$ 作用，便定义出所要的\trPinStructure{}$P^{\mathrm{fix},\#}$。
现在设 $(L_0^{\#},L_1^{\#})$ 是两个这样的带\trBraneStructure{}的\trLagrange\trSubmanifold，而 $(H,J)$ 是如\cref{thm:floer-datum-assumption}中的一个\trFloerDatum{}。如前面所提到的，把它限制到\trFixedPoint{}集上，就给出 $(L_0^{\mathrm{fix}},L_1^{\mathrm{fix}})$ 的一个\trFloerDatum{}。取某个含于 $M^{\mathrm{fix}}$ 之内的 $y\in\scr{C}(L_0,L_1)$，并把同一\trObject{}视为 $\scr{C}(L_0^{\mathrm{fix}},L_1^{\mathrm{fix}})$ 中的元素时记作 $y^{\mathrm{fix}}$。那么，配上如上定义的\trFixedPoint{}集上的诱导\trBraneStructure{}，便有
\begin{lemma}\label{thm:equivariant-index-orientation}
$i(y^{\mathrm{fix}})=i(y)$，并且 $o(y^{\mathrm{fix}})\cong o(y)$。
\end{lemma}
\begin{proof}
依循 \eqref{eq:orientation-operator-index}，我们把 $i(y)$ 与 $o(y)$ 设想为一个\trOrientationOperator{}的\trIndex{}与\trDeterminantLine，这个算子关联于由诸 $L_k^{\#}$ 在 $y$ 处的\trTangentSpace{}给出的那对\trLinear\trLagrange\trBrane{}$\Lambda^{\#}_{k,y}$。所要的断言是：可以把这个\trOrientationOperator{}选得成为 $\ZZ/2$-\trEquivariant{}的，其不变部分是 $y^{\mathrm{fix}}$ 的\trOrientationOperator{}，而反不变部分\trInvertible{}。这类似于\cref{thm:regularity-fixed-set}与\cref{thm:regularity-polygons}的证明。为使记号简单，我们只对 $H=0$ 的情形写出细节。此时 $y$ 是 $L_0$ 与 $L_1$ 的一个不变交点，而我们可以如下构造一个\trOrientationOperator{}$D_{H,y}$。取\trVectorBundle{}$E=H\times TM_y$，配上一个在无穷远处由我们的\trFloerDatum{}确定的非常值\trComplexStructure{}$J_E$（这意味着：若 $\epsilon_\zeta$ 是 $H$ 的那一条\trStriplikeEnd{}，则当 $s\ll 0$ 时 $J_E(\epsilon_\zeta(s,t))=J_t$）。我们给 $E$ 配上平凡\trConnection{}$\nabla$。\trSubbundle{}$F=F_\lambda\subset E|\partial H$ 由 $Gr(TM_y)$ 中一条从 $\lambda(0)$ 连到 $\lambda(1)$ 的路径 $\lambda$ 给出。最后，为使这成为一个真正的\trOrientationOperator{}，还需要为 $\lambda$ 选一个\trBraneStructure{}$(\alpha_\lambda^{\#},P_\lambda^{\#})$，这当然会对该路径的\trHomotopy{}类施加限制。
在这个类之内，可以把 $J_E$ 选得在 $D\iota_M|TM_y$ 下不变，并让路径 $\lambda$ 由不变的\trLagrange{}子空间组成，于是 $D_{H,y}=\bar\partial_\nabla$ 就成为 $\ZZ/2$-\trEquivariant{}的。此外，$\lambda$ 的反不变部分是 $\nu_y\cong\CC$ 中的一族直线 $\lambda^-(s)$，可以取其形为 $\lambda^-(s)=e^{j\sigma(s)}\RR$，其中某个函数 $\sigma\colon[0;1]\to\RR$ 从 $\sigma(0)=\sigma_0$ 走到 $\sigma(1)=\sigma_1$。在此之下，不变部分按照上文对\trLagrange\trSubmanifold{}的同一套做法继承一个诱导\trBraneStructure{}，这使得 $D^+_{H,y}$ 成为 $y^{\mathrm{fix}}$ 的一个\trOrientationOperator{}。另一方面，条件 \eqref{eq:normal-angle-chamber}——它在我们的情形恰好是说 $\sigma_1-\sigma_0\in(-\pi;0)$——蕴含 $\mathrm{index}(D^-_{H,y})=0$。再由\cref{thm:winding-injectivity}，即知 $D^-_{H,y}$ 是\trInvertible{}的。
\end{proof}
\begin{remark}\label{thm:index-free-signs-rephrase}
或许有意思的是，把这个论证换成一种避开\trIndex{}理论的说法来重述，即改为依赖 \eqref{eq:absolute-index-orientation-space} 中\trIndex{}与\trOrientation{}空间的初等定义。具体地说，选 $\lambda$ 使得它在上述条件之外还落在 $\scr{P}^-\,Gr(V)$ 中。此时，\trIndex{}与\trDeterminantLine{}便只依赖于诸交点 $\lambda(s)\cap\lambda(1)$。然而，我们可以把函数 $\sigma$ 选成严格递减的，于是 $\sigma_1-\sigma_0\in(-\pi;0)$ 这一事实就意味着：对一切 $s<1$ 都有 $\lambda^-(s)\cap\lambda^-(1)=0$。因此，所有交点都集中在不变部分里。
\end{remark}
设我们的\trFloerDatum{}是\trRegular{}的。取一个落在 $M^{\mathrm{fix}}$ 内、是该\trModuliSpace{}的\trIsolated{}点的 $u\in\scr{M}^*_Z(y_0,y_1)$。一如\cref{thm:regularity-fixed-set}的证明，我们把同一映射视为 $\scr{M}^*_Z(y_0^{\mathrm{fix}},y_1^{\mathrm{fix}})$ 中的元素时记作 $u^{\mathrm{fix}}$。我们接下来断言：相对于来自 \cref{thm:equivariant-index-orientation} 的那些\trIsomorphism，$u$ 与 $u^{\mathrm{fix}}$ 对各\trFloerComplex{}的\trDifferential{}的贡献具有相同的符号。用\cref{sub:determining-grading-and-signs}的记号来说，这意味着下图交换：
\begin{equation}\label{eq:equivariant-cu-diagram}
\begin{tikzcd}
o_{y_1}\ar[r,"c_u"]\ar[d,"\cong"] & o_{y_0}\ar[d,"\cong"] \\
o_{y_1^{\mathrm{fix}}}\ar[r,"c_{u^{\mathrm{fix}}}"] & o_{y_0^{\mathrm{fix}}}.
\end{tikzcd}
\end{equation}
要看出这一点，需要回想 \eqref{eq:disc-dim-orientation}：$c_u$ 是这样得到的——把\trLinearization{}算子 $D_{Z,u}$ 与\trOrientationOperator{}$D_{H,y_1}$ \trGluing{}起来，再把所得结果\trDeformation{}到 $D_{H,y_0}$。在我们的情形，所有这些算子都是 $\ZZ/2$-\trEquivariant{}的；它们的不变部分可以等同于 $M^{\mathrm{fix}}$ 中相应的算子（后者给出 $c_{u^{\mathrm{fix}}}$），而它们的反不变部分是\trInvertible{}的，从而对\trDeterminantLine{}的贡献是平凡的。此外，从 $D_{Z,u}\#D_{H,y_1}$ 到 $D_{H,y_0}$ 的\trHomotopy{}全程经由同形的算子，这意味着反不变部分始终保持\trInvertible{}。同样的论证也适用于\cref{thm:regularity-polygons}的情形中那些不变的\trIsolated{}\trPseudoholomorphic\trPolygon{}的贡献 $c_{r,u}$。
